Phần này tạo bộ não của hệ thống: một hàm Lambda xử lý toàn bộ nghiệp
vụ, và một API Gateway làm cửa vào.

\textbf{Lưu ý bảo mật (Auth):} Trong phạm vi workshop này, logic băm mật
khẩu và quản lý người dùng được lập trình trực tiếp trong API để bạn dễ
hình dung luồng dữ liệu. Trên môi trường Production thực tế, bạn nên ủy
quyền toàn bộ chức năng này cho \textbf{Amazon Cognito} để đảm bảo tiêu
chuẩn bảo mật.

Thứ tự làm quan trọng: \textbf{IAM role trước, Lambda sau, API Gateway
cuối}. Vì Lambda cần role lúc tạo, và API Gateway cần Lambda đã tồn tại.

\vspace{0.5em}\noindent\textbf{Bước 1. Tạo IAM role theo nguyên tắc quyền tối
thiểu}\par

Lambda cần quyền đọc ghi DynamoDB và gọi Personalize. Cách nhanh nhất là
gán \texttt{AdministratorAccess}, nhưng đó là thói quen xấu: nếu mã
nguồn có lỗ hổng, kẻ tấn công chiếm được toàn bộ tài khoản.

Ta chỉ cấp đúng thứ cần.

\begin{enumerate}
\def\labelenumi{\arabic{enumi}.}
\tightlist
\item
  Mở \href{https://console.aws.amazon.com/iam/home\#/roles}{IAM
  console}, chọn \textbf{Roles} rồi \textbf{Create role}
\item
  \textbf{Trusted entity type}: AWS service. \textbf{Use case}: Lambda.
  Chọn \textbf{Next}
\item
  Tìm và chọn \texttt{AWSLambdaBasicExecutionRole} --- quyền ghi log vào
  CloudWatch
\item
  Chọn \textbf{Next}, đặt tên \texttt{fcj-recsys-lambda-role}, chọn
  \textbf{Create role}
\end{enumerate}

Giờ thêm quyền riêng cho DynamoDB và Personalize:

\begin{enumerate}
\def\labelenumi{\arabic{enumi}.}
\setcounter{enumi}{4}
\tightlist
\item
  Mở role vừa tạo, chọn \textbf{Add permissions} rồi \textbf{Create
  inline policy}
\item
  Chuyển sang tab \textbf{JSON}, dán nội dung sau (thay
  \texttt{\textless{}ACCOUNT-ID\textgreater{}} bằng số tài khoản của
  bạn):
\end{enumerate}

\begin{Shaded}
\begin{Highlighting}[]
\FunctionTok{\{}
  \DataTypeTok{"Version"}\FunctionTok{:} \StringTok{"2012{-}10{-}17"}\FunctionTok{,}
  \DataTypeTok{"Statement"}\FunctionTok{:} \OtherTok{[}
    \FunctionTok{\{}
      \DataTypeTok{"Effect"}\FunctionTok{:} \StringTok{"Allow"}\FunctionTok{,}
      \DataTypeTok{"Action"}\FunctionTok{:} \OtherTok{[}
        \StringTok{"dynamodb:GetItem"}\OtherTok{,} \StringTok{"dynamodb:PutItem"}\OtherTok{,} \StringTok{"dynamodb:UpdateItem"}\OtherTok{,}
        \StringTok{"dynamodb:DeleteItem"}\OtherTok{,} \StringTok{"dynamodb:Query"}\OtherTok{,} \StringTok{"dynamodb:Scan"}\OtherTok{,}
        \StringTok{"dynamodb:BatchGetItem"}\OtherTok{,} \StringTok{"dynamodb:BatchWriteItem"}
      \OtherTok{]}\FunctionTok{,}
      \DataTypeTok{"Resource"}\FunctionTok{:} \StringTok{"arn:aws:dynamodb:ap{-}southeast{-}1:\textless{}ACCOUNT{-}ID\textgreater{}:table/*"}
    \FunctionTok{\}}\OtherTok{,}
    \FunctionTok{\{}
      \DataTypeTok{"Effect"}\FunctionTok{:} \StringTok{"Allow"}\FunctionTok{,}
      \DataTypeTok{"Action"}\FunctionTok{:} \StringTok{"personalize:GetRecommendations"}\FunctionTok{,}
      \DataTypeTok{"Resource"}\FunctionTok{:} \StringTok{"*"}
    \FunctionTok{\}}
  \OtherTok{]}
\FunctionTok{\}}
\end{Highlighting}
\end{Shaded}

\begin{enumerate}
\def\labelenumi{\arabic{enumi}.}
\setcounter{enumi}{6}
\tightlist
\item
  Đặt tên policy \texttt{fcj-recsys-lambda-policy}, chọn \textbf{Create
  policy}
\end{enumerate}

Lấy Account ID bằng lệnh
\texttt{aws\ sts\ get-caller-identity\ -\/-query\ Account\ -\/-output\ text}

\vspace{0.5em}\noindent\textbf{Bước 2. Đóng gói mã
nguồn}\par

Lambda cần một file zip chứa mã nguồn và thư viện phụ thuộc:

\begin{Shaded}
\begin{Highlighting}[]
\BuiltInTok{cd}\NormalTok{ backend}
\ExtensionTok{npm}\NormalTok{ install }\AttributeTok{{-}{-}production}
\FunctionTok{zip} \AttributeTok{{-}r}\NormalTok{ ../function.zip index.mjs node\_modules}
\end{Highlighting}
\end{Shaded}

\vspace{0.5em}\noindent\textbf{Bước 3. Tạo hàm
Lambda}\par

\begin{enumerate}
\def\labelenumi{\arabic{enumi}.}
\tightlist
\item
  Mở
  \href{https://ap-southeast-1.console.aws.amazon.com/lambda/home?region=ap-southeast-1}{Lambda
  console}
\item
  Chọn \textbf{Create function}
\item
  Điền:

  \begin{itemize}
  \tightlist
  \item
    \textbf{Function name}: \texttt{fcj-recsys-api}
  \item
    \textbf{Runtime}: \textbf{Node.js 20.x}
  \item
    \textbf{Architecture}: x86\_64
  \item
    \textbf{Execution role}: chọn \textbf{Use an existing role}, chọn
    \texttt{fcj-recsys-lambda-role}
  \end{itemize}
\item
  Chọn \textbf{Create function}
\item
  Ở tab \textbf{Code}, chọn \textbf{Upload from} rồi \textbf{.zip file},
  tải \texttt{function.zip} lên
\item
  Vào \textbf{Runtime settings}, chọn \textbf{Edit}, đặt
  \textbf{Handler} thành \texttt{index.handler}
\end{enumerate}

\begin{figure}
\centering
\pandocbounded{\includegraphics[keepaspectratio,alt={Tạo Lambda}]{5-Workshop/5.5/create-lambda.png}}
\caption{Tạo Lambda}
\end{figure}

\vspace{0.5em}\noindent\textbf{Bước 4. Tăng
timeout}\par

Mặc định Lambda dừng sau 3 giây, không đủ khi phải gọi Personalize.

\begin{enumerate}
\def\labelenumi{\arabic{enumi}.}
\tightlist
\item
  Mở tab \textbf{Configuration}, chọn \textbf{General configuration},
  chọn \textbf{Edit}
\item
  Đặt \textbf{Timeout} là \texttt{30} giây, \textbf{Memory} là
  \texttt{512} MB
\item
  Chọn \textbf{Save}
\end{enumerate}

\vspace{0.5em}\noindent\textbf{Bước 5. Tạo HTTP
API}\par

\begin{enumerate}
\def\labelenumi{\arabic{enumi}.}
\tightlist
\item
  Mở
  \href{https://ap-southeast-1.console.aws.amazon.com/apigateway/main/apis?region=ap-southeast-1}{API
  Gateway console}
\item
  Ở ô \textbf{HTTP API}, chọn \textbf{Build}
\item
  Chọn \textbf{Add integration}, kiểu \textbf{Lambda}, chọn hàm
  \texttt{fcj-recsys-api}
\item
  \textbf{API name}: \texttt{fcj-recsys-api}, chọn \textbf{Next}
\item
  Ở phần \textbf{Configure routes}, khai báo một route bắt tất cả:

  \begin{itemize}
  \tightlist
  \item
    \textbf{Method}: \texttt{ANY}
  \item
    \textbf{Resource path}: \texttt{/\{proxy+\}}
  \item
    \textbf{Integration target}: hàm Lambda vừa chọn
  \end{itemize}
\item
  Chọn \textbf{Next}, giữ stage \texttt{\$default} với
  \textbf{Auto-deploy} bật, chọn \textbf{Create}
\end{enumerate}

Route \texttt{/\{proxy+\}} chuyển mọi đường dẫn về cùng một Lambda, để
mã nguồn tự định tuyến bên trong. Cách này giữ cấu hình API Gateway đơn
giản, đổi lại việc định tuyến nằm trong mã. Với 22 route thì đây là đánh
đổi hợp lý.

\vspace{0.5em}\noindent\textbf{Bước 6. Bật CORS}\par

Giao diện chạy trên tên miền CloudFront, API nằm ở tên miền khác. Trình
duyệt sẽ chặn nếu không khai báo CORS.

\begin{enumerate}
\def\labelenumi{\arabic{enumi}.}
\tightlist
\item
  Trong API vừa tạo, chọn \textbf{CORS} ở menu bên trái, chọn
  \textbf{Configure}
\item
  Khai báo:

  \begin{itemize}
  \tightlist
  \item
    \textbf{Access-Control-Allow-Origin}: \texttt{*}
  \item
    \textbf{Access-Control-Allow-Headers}:
    \texttt{content-type,\ authorization}
  \item
    \textbf{Access-Control-Allow-Methods}:
    \texttt{GET,\ POST,\ PUT,\ DELETE,\ OPTIONS}
  \end{itemize}
\item
  Chọn \textbf{Save}
\end{enumerate}

Lỗi CORS chỉ hiện trong console trình duyệt, còn thử bằng \texttt{curl}
thì vẫn chạy bình thường. Nếu giao diện báo không kết nối được máy chủ
trong khi \texttt{curl} vẫn ổn, hãy bấm F12 xem tab Console trước tiên.

\vspace{0.5em}\noindent\textbf{Bước 7. Nối giao diện với
API}\par

Copy \textbf{Invoke URL} ở trang chi tiết API, rồi quay lại thư mục giao
diện:

\begin{Shaded}
\begin{Highlighting}[]
\BuiltInTok{cd}\NormalTok{ frontend}
\BuiltInTok{echo} \StringTok{"VITE\_API\_BASE\_URL=https://xxxxx.execute{-}api.ap{-}southeast{-}1.amazonaws.com"} \OperatorTok{\textgreater{}}\NormalTok{ .env}
\ExtensionTok{npm}\NormalTok{ run build}
\ExtensionTok{aws}\NormalTok{ s3 sync dist/ s3://fcj{-}recsys{-}frontend{-}}\OperatorTok{\textless{}}\NormalTok{tên{-}bạn}\OperatorTok{\textgreater{}}\NormalTok{/ }\AttributeTok{{-}{-}delete}
\ExtensionTok{aws}\NormalTok{ cloudfront create{-}invalidation }\AttributeTok{{-}{-}distribution{-}id} \OperatorTok{\textless{}}\NormalTok{ID{-}CUA{-}BAN}\OperatorTok{\textgreater{}}\NormalTok{ {-}{-}paths }\StringTok{"/*"}
\end{Highlighting}
\end{Shaded}

\vspace{0.5em}\noindent\textbf{Kiểm tra}\par

\begin{Shaded}
\begin{Highlighting}[]
\ExtensionTok{curl}\NormalTok{ https://xxxxx.execute{-}api.ap{-}southeast{-}1.amazonaws.com/products}
\end{Highlighting}
\end{Shaded}

Trả về danh sách sản phẩm dạng JSON là thành công. Mở lại trang
CloudFront, sản phẩm phải hiển thị.
